%%
%% fall-lecture22.tex
%% 
%% Made by Alex Nelson <pqnelson@gmail.com>
%% Login   <alex@lisp>
%% 
%% Started on  2026-11-19T10:59:03-0800
%% Last update 2026-11-19T10:59:03-0800
%% 

\lecture{}

\begin{node}
We were in the middle of proving Property 3: $c_{1}(\gamma_{1}^{1})\in H^{2}(\CP^{1};\ZZ)\iso\ZZ$
s the generator of $\ZZ$. So $c_{1}(\gamma_{1}^{1})=\pm e((\gamma_{1}^{1})_{\RR})$.
First, we need to review the infinite complex projective space $\CP^{\infty}$.
\end{node}

\begin{definition}
A \define{Complex $n$-Frame} in $\CC^{n+k}$ is a $n$-tuple of linearly
independent [over $\CC$] vectors living in $\CC^{n+k}$.
\end{definition}

\begin{definition}
The \define{Complex Stiefel Manifold} is the set $V_{n}(\CC^{n+k})$ of
all complex $n$-frames, which is a subset of $(\CC^{n+k})^{n}$ and we
equip it with the subspace topology.
\end{definition}

\begin{definition}
The \define{Complex Grassmannian Manifold} $G_{n}(\CC^{n+k})$ is the
set of all complex linear $n$-subspaces $\CC^{n+k}$. It has a topology
given by the quotient topology using the quotient map
\begin{align}
q\colon & V_{n}(\CC^{n+k})\to G_{n}(\CC^{n+k})\\
 & (z_{1},\dots,z_{n})\mapsto \Span_{\CC}(z_{1},\dots,z_{n})\nonumber.
\end{align}
\end{definition}

\begin{remark}
Observe $G_{n}(\CC^{n+k})$ is a complex manifold and has [complex]
dimension $nk$. All we really need to check is that it has a complex structure.
\end{remark}

\begin{remark}
When $n=1$, $G_{1}(\CC^{1+k})\iso\CP^{k}$ homeomorphic (diffeomorphic, even!)
where $\CP^{k}$ is the set of all complex lines in $\CC^{k+1}$. Note
$G_{1}(\CC^{2})\iso\CP^{1}\iso\sphere{2}$ when $n=k=1$.
\end{remark}

\begin{definition}
Denote $C^{\infty}=\{(z_{1},z_{2},\dots)\mid\mbox{only finitely many entries are nonzero}\}$
(possibly equipped with the directed limit topology).
Then we can take the directed limit
\begin{equation}
G_{n}(\CC^{\infty}):=\varinjlim_{k\to\infty}G_{n}(\CC^{n+k})
\end{equation}
which is called the \define{Infinite Complex Grassmann Manifold}
equipped with the directed limit topology.
\end{definition}

\begin{definition}
The \define{Universal Complex Bundle} over $G_{n}(\CC^{\infty})$
is $\gamma^{n}:=\{(X,v)\in G_{n}(\CC^{\infty})\times\CC^{\infty}\mid v\in X\}$.
\end{definition}

\begin{theorem}
The cohomology ring $H^{*}(G_{n}(\CC^{\infty});\ZZ)\iso\ZZ[c_{1}(\gamma^{n}),\dots,c_{n}(\gamma^{n})]$
where $\gamma^{n}$ is the universal complex bundle over $G_{n}(\CC^{\infty})$.
\end{theorem}

\begin{remark}
The analogous real version uses $w_{i}(\gamma^{n})$ instead of $c_{i}(\gamma^{n})$.
\end{remark}

(Proof for $n=1$ case I will not transcribe, and the $n>1$ case
omitted because it is tedious and not very enlightening.)

\begin{fact}
The cohomology ring for the infinite-dimensional Grassmann manifold is
$H^{*}(G_{n}(\CC^{\infty});\ZZ)\iso\ZZ[c_{1}(\gamma^{n}),\dots,c_{n}(\gamma^{n})]$.
\end{fact}

\begin{lemma}
Let $m,n\in\NN$.
There exists a unique polynomial $P_{m,n}(c_{1},\dots,c_{m},c'_{1},\dots,c'_{n})$
in $m+n$ unknowns with integer coefficients such that for 
every $\CC^{m}$-bundle $\omega$ and $\CC^{n}$-bundle $\phi$ both over
the same paracompact space $B$, we have
\begin{equation}
c(\omega\oplus\phi)=P_{m,n}(c_{1}(\omega),\dots,c_{m}(\omega),c_{1}(\phi),\dots,c_{n}(\phi)).
\end{equation}
This polynomial $P_{m,n}$ is called the \define{Universal Polynomial}.
\end{lemma}

We will prove $P_{m,n}=(1+c_{1}(\omega)+\dots+c_{m}(\omega))(1+c_{1}(\phi)+\dots+c_{n}(\phi))$.